\documentclass[11pt]{article}
\usepackage[margin=1in]{geometry}
\usepackage{booktabs}
\usepackage{amsmath}
\usepackage{hyperref}
\usepackage{microtype}

\title{Boundedness, Not the Posterior: A Preregistered Replication of
TreeOPSD on Qwen3-4B}
\author{Stephen J. Ronan\\RonanLabs\\Blackhawk, CA, USA}
\date{October 2026}

\begin{document}
\maketitle

\begin{abstract}
We preregister and run a replication of TreeOPSD (\textit{Tree-Aware
Posterior Distillation for Long-Horizon Agentic Reinforcement Learning},
arXiv:2610.23181) on Qwen3-4B-Thinking-2507 with search-agent training.
Against the paper's own control, unclipped TreePO, TreeOPSD passes both of
our preregistered bars at three seeds: pooled exact-match gaps of $+3.76$
$[+2.81, +4.74]$ on a four-set multi-hop QA average and $+3.45$
$[+2.13, +4.78]$ on the long-horizon BC+ benchmark (830 questions). Both
passes are fragile: seed-level intervals at $n=3$ include zero, and the
largest contributor is one control seed that began to destabilise at the
paper's dose. Against an added control, TreePO with ordinary ratio
clipping, the BC+ gap is $+0.08$ $[-1.24, +1.41]$ and the QA gap is
$+1.98$ $[+1.03, +2.94]$ pooled but declining across seeds
($3.78, 2.10, 0.05$). The posterior term changes the per-update target by
a median $1.05\%$ at the paper's dose; a preregistered dose sweep that
raises this drift sevenfold produces no growing accuracy gap, and at
$4\times$ dose the unclipped control collapses in training in two of two
seeds while both bounded variants train through. The budget match claimed
for tree sampling holds in generated tokens ($1.03\times$ GRPO) but not in
trained tokens ($1.48\times$). Best-checkpoint selection shrinks the gap
toward zero. We conclude that the supported mechanism in this setting is
update boundedness, not posterior credit assignment, and that bounded
tree backup itself is an effective and, once clipped, low-variance
($0.80$ EM seed sd) method for training a 4B search agent.
\end{abstract}

\section{Introduction}

TreeOPSD \cite{treeopsd} trains long-horizon agents by sampling a tree of
continuations from shared prefixes, backing terminal rewards up the tree
into turn-level values, converting those values into a bounded teacher
distribution relative to a frozen behaviour policy, and matching the
student to the teacher with a forward-KL self-distillation loss. The
method is critic-free and is presented as a solution to credit assignment
under sparse terminal rewards. Reported effects on multi-hop QA for
Qwen3-4B-Thinking-2507 are small (TreeOPSD $44.13$ vs.\ TreePO $43.40$
vs.\ GRPO $43.37$ MHQA average, five-seed means), with larger effects
($+2.9$ to $+4.0$) on a deep-research track.

A replication of this claim faces a confound in the original design. The
paper's tree control, TreePO, applies its backed-up values without any
ratio clip or bound, whereas TreeOPSD's teacher is bounded by
construction. Outperformance against that control is therefore consistent
with two mechanisms: (i) the posterior teacher supplies turn-level credit
signal, as claimed; or (ii) any bound on the update improves stability
and downstream accuracy. We resolve the confound by adding a control arm
that keeps TreePO's sampler and backup but applies ordinary ratio
clipping. We preregistered all arms, metrics, checkpoints and
pass/kill bars before training (Section~\ref{sec:design}); the
preregistration (\texttt{PREREG-EXP1.md}, one pre-training amendment) ships
with our code and result packages.

\paragraph{Contributions.} (1) A five-arm, three-seed replication with
frozen retrieval and a fixed evaluation checkpoint. (2) A direct
measurement of the posterior term's per-update effect (M1) and a
preregistered dose--response test of its causal role. (3) Evidence that
unclipped tree backup is systematically unstable (collapse in $2/2$ seeds
at $4\times$ dose; destabilisation in $1/3$ seeds at $1\times$), and that
either bound prevents it. (4) A token-accounting audit showing that the
tree budget match holds for generated tokens ($1.03\times$) but not for
trained tokens ($1.48\times$).

\section{Design}
\label{sec:design}

\paragraph{Model and task.} Qwen3-4B-Thinking-2507, ReAct-style search
agents with a frozen retrieval backend, trained for 100 steps
\mbox{($M,N)=(4,2)$} trees for the tree arms) and evaluated greedily,
one sample per question, at the fixed step-100 checkpoint. Tier~1 is the
unweighted mean exact match (EM) over HotpotQA dev (7{,}405 questions),
2WikiMultihopQA, MuSiQue and Bamboogle (125). Tier~2 is BC+ (830
questions), our closed-corpus long-horizon track and the preregistered
headline. Training and full evaluation ran on two DGX Station GB300
systems and an RTX PRO 6000 sidecar; all C-vs-E scores reported as
headline were re-scored on GB300 hardware so the deciding comparison
shares one device class (as-scored values differ by at most $0.49$ EM and
change no conclusion).

\paragraph{Arms.} A: GRPO, $G=8$ (seed 0). B: TreePO literal, the paper's
control, likelihood-$\omega$ backup, prompt-normalised, no ratio or clip
(seeds 0--2). C: TreeOPSD (seeds 0--2). D: GRPO, $G=12$, the
trained-token-matched GRPO (seed 0). E: TreePO with ordinary ratio
clipping (seeds 0--2). All tree arms share the sampler, tree shape, data
and retrieval.

\paragraph{Statistics.} Rows are paired across runs by prompt hash.
Gaps carry paired bootstrap intervals (10{,}000 resamples over questions,
stratified by benchmark; pooled gaps average the per-question difference
over seeds before resampling). We also report seed-level means, standard
deviations and $t$ intervals ($n=3$). The bootstrap intervals carry
evaluation noise only; the seed intervals carry training-seed noise only.
We report both throughout, because in this setting they disagree, which
is itself a result.

\paragraph{Preregistered bars.} F1 (posterior does the work): kill if
$|\text{TreeOPSD}-\text{TreePO}|<0.25$ with the interval excluding the
paper's $+0.73$ on tier~1, or if the gap is $<1.0$ on tier~2 (where the
paper claims $+2.90$); pass at $\geq +0.5$ (tier~1) / $\geq +1.5$
(tier~2) with intervals excluding zero. Mechanism check M1: let
$\mathrm{drift\_rel}=\overline{|\Gamma-V|}/\overline{|V|}$ over inner-loop
blocks 2--4, where $\Gamma$ and $V$ are the TreeOPSD and TreePO targets.
If $\mathrm{drift\_rel}<0.02$, the posterior term is declared inactive at
that dose whatever the EM shows, and a preregistered dose sweep
(EXP-1c) follows. F2 (budget): kill if any of generated tokens, tool
calls or GPU-hours exceeds GRPO by $>1.15\times$, or if budget matching
removes $>50\%$ of the gap. F3 (selection): kill if the gap's interval
includes zero at the fixed checkpoint but excludes zero under
best-in-window selection. F4 (retrieval): frozen by construction. F5
(horizon): gain should grow with hop count; our MHQA design caps at five
turns, so F5 evidence here is weak by construction and we claim no trend.

\section{Results}

\begin{table}[t]
\centering
\caption{Step-100 exact match. Tier~1 (QA avg) is the mean over the four
QA sets; Bamboogle (125 questions) is included in the average but its
column is omitted for width. C is shown GB300-scored (see text). A and D ran seed 0 only. The
untrained base scores $30.8$ on frozen HotpotQA dev (paper ReAct $31.61$)
and $22.8/6.5$ on the 500/200-question interim subsets.}
\label{tab:arms}
\begin{tabular}{lccccc}
\toprule
Run & BC+ & HotpotQA & 2Wiki & MuSiQue & QA avg \\
\midrule
A GRPO G=8 s0 & 10.00 & 43.44 & 48.55 & 18.41 & 39.80 \\
B TreePO s0/s1/s2 & 10.48/7.23/6.27 & 42.58/45.97/39.54 & 49.93/53.78/41.57 & 18.57/20.61/14.05 & 39.17/42.89/34.99 \\
C TreeOPSD s0/s1/s2 & 10.48/9.04/13.61 & 46.71/46.29/45.43 & 51.67/50.65/48.74 & 22.35/21.10/21.27 & 43.58/42.71/41.46 \\
D GRPO G=12 s0 & 7.83 & 44.09 & 51.19 & 20.32 & 41.50 \\
E TreePO-clipped s0/s1/s2 & 12.29/12.05/8.55 & 44.74/45.12/44.11 & 47.73/48.24/49.48 & 20.36/19.49/18.45 & 39.81/40.61/41.41 \\
\bottomrule
\end{tabular}
\end{table}

\paragraph{Posterior vs.\ bounded (C$-$E): the deciding comparison.}
Table~\ref{tab:gaps} gives the hardware-matched result. On tier~1 QA the
pooled gap is $+1.98$ $[+1.03, +2.94]$, meeting the script's pass bar;
per-seed it declines $3.78, 2.10, 0.05$, and the seed-level interval
$[-2.66, +6.61]$ includes zero. On tier~2 BC+, where the paper's large
effect is claimed, the pooled gap is $+0.08$ $[-1.24, +1.41]$ (per-seed
$-1.81, -3.01, +5.06$), meeting the kill bar. The clipped control's tier~1
seed sd ($0.80$) is the smallest of any arm (C $1.07$, B $3.95$). We read
tier~1 as a real but seed-fragile pooled edge of at most $\approx\!2$ EM
whose cause Section~\ref{sec:mech} does not allow us to credit to the
posterior mechanism, and tier~2 as no posterior-specific effect.

\paragraph{Against the paper's control (C$-$B).} Pooled gaps are tier~1
$+3.76$ $[+2.81, +4.74]$ and tier~2 $+3.45$ $[+2.13, +4.78]$: passes by the
letter of F1's EM bars. Seed-level intervals ($[-4.06, +11.57]$ and
$[-6.13, +13.03]$) include zero. B seed~2 began to destabilise in its
final 10 steps at the paper's dose: reward fell from $0.54$--$0.70$ to
$0.22$--$0.40$, the length-clip ratio rose from $0.00$ to $0.27$, and the
checkpoint issued $3.52$ tool calls per BC+ question against $1.81/1.97$
for B s0/s1. Excluding B s2, the two-seed gaps are $+2.40$ (tier~1) and
$+1.51$ (tier~2). Clipping alone (E$-$B) improves BC+ in all three seeds
($+1.81, +4.82, +2.29$).

\begin{table}[t]
\centering
\caption{Pooled gaps (question bootstrap, 10{,}000 resamples) with
per-seed point values and seed-level $t$ intervals ($n=3$). C$-$E is
hardware-matched; as-scored it reads $+2.17\,[+1.29,{+}3.06]$ (tier~1)
and $+0.48\,[-0.84,{+}1.81]$ (tier~2).}
\label{tab:gaps}
\begin{tabular}{lllll}
\toprule
Gap & Tier & Pooled [95\% CI] & Per-seed & Seed $t$ interval \\
\midrule
C$-$E & QA & $+1.98$ $[+1.03, +2.94]$ & $3.78, 2.10, 0.05$ & $[-2.66, +6.61]$ \\
C$-$E & BC+ & $+0.08$ $[-1.24, +1.41]$ & $-1.81, -3.01, 5.06$ & $[-10.74, +10.90]$ \\
C$-$B & QA & $+3.76$ $[+2.81, +4.74]$ & $4.50, 0.31, 6.47$ & $[-4.06, +11.57]$ \\
C$-$B & BC+ & $+3.45$ $[+2.13, +4.78]$ & $-0.36, 3.37, 7.35$ & $[-6.13, +13.03]$ \\
C$-$A (s0) & QA / BC+ & $+3.78$ $[+2.41, +5.21]$ / $+0.48$ $[-1.69, +2.65]$ & -- & -- \\
C$-$D (s0) & QA / BC+ & $+2.08$ $[+0.69, +3.53]$ / $+2.65$ $[+0.60, +4.82]$ & -- & -- \\
\bottomrule
\end{tabular}
\end{table}

\section{Mechanism: M1 and the dose response}
\label{sec:mech}

At the paper's dose ($E=4$ inner-loop blocks per step) M1 trips:
$\mathrm{drift\_rel}$ median $1.05\%$ over the full run and $1.56\%$ in
the last quarter, below the preregistered $2\%$ inactivity line. Backup
weights are near-one-hot in both tree arms
($\omega_{\mathrm{peaked}} \approx 0.85$--$0.91$), so the backup is
effectively single-child in this setting. EXP-1c raised the dose. At
$E=8$, drift rose (median $1.23\%$, last-quarter $1.89\%$) and the tier~1
subset gap was $-0.8$ EM: no growing accuracy gap. At $E=16$ (drift
$4.79\%/7.74\%$, $34\%$ of calls above $10\%$) unclipped TreePO collapsed
in training: reward $0.002$ and length-clip ratio $0.99$ by step~33; the
step-40 checkpoint scores $0.0/0.0$ on the interim subsets. A second seed
collapsed the same way (subset scores $17.8/0.0$; BC+ fully collapsed, QA
degraded but answering). At the same dose, clipped TreePO trained through
(reward $0.49$--$0.64$, clip ratio $\leq 0.10$; subsets $34.2/8.5$) as did
TreeOPSD ($0.50$--$0.69$; $35.4/5.0$). Boundedness, in either form, is what
separates survival from collapse; the posterior term adds no measurable
signal beyond it at any tested dose. Combined with B s2's destabilisation
at $E=4$, unclipped tree backup failed fully in $2/2$ seeds at $4\times$
dose and partially in $1/3$ seeds at $1\times$; neither bounded variant
failed in any run.

\paragraph{Budget (F2).} Tree sampling is $1.03\times$ GRPO $G{=}8$ in
generated tokens per prompt and $1.48\times$ in trained tokens per prompt
($41{,}064$ vs.\ $27{,}779$ in the matched smoke runs), because shared
prefixes are re-scored per branch during training. Token-matched GRPO
(arm D) removes $44\%$ of the tier~1 C$-$A gap, below the $50\%$ F2b kill
line; on BC+ the $G{=}8$ gap is $\approx 0$ and the ratio is undefined.
The claimed budget match therefore depends on accounting in generated
tokens only.

\paragraph{Selection (F3).} Best-in-window scoring over checkpoints
10/25/50/75/100 on fixed subsets shrinks C$-$B toward zero (per-seed mean
$+0.3$ tier~1 / $+0.5$ tier~2, vs.\ $+1.5/+2.5$ at the final checkpoint
on the same subsets), chiefly by rescuing B s2's pre-destabilisation peak
($36.6$ at step~75 vs.\ $32.8$ at step~100, qa500i). The F3 kill pattern
is not met and the shift runs in the opposite direction; selection favours
the control. Our grid is coarser than the preregistration's every-5-steps
ideal, which is conservative against this conclusion. Retrieval was
frozen by construction (F4); an eval-drift control (B s0 Bamboogle
re-scored after the study) returned an identical $45.6$.

\section{Discussion and limitations}

Three claims survive this replication. Bounded tree backup trains a 4B
search agent effectively (base $22.8$ EM on our 500-question QA subset;
$40.61 \pm 0.80$ and $42.58 \pm 1.07$ full-set QA means for arms E and C).
Unclipped tree backup is fragile, systematically so at high dose.
TreeOPSD outperforms the control as the paper defines it. The claim that
does not survive is the mechanism: the posterior teacher's per-update
effect is $\approx\!1\%$, its dose response is flat, and against a bounded
control its advantage vanishes on the long-horizon track. The residual
tier~1 C$-$E pooled edge is consistent with second-order dynamics we
observed but did not isolate (C sustained $\approx\!43$K trained tokens
per prompt where B's fell to $\approx\!33$K, and learned a longer,
search-heavier style: $14.2$ vs.\ $10.6$ tool calls per prompt late in
training) and with seed noise; $n=3$ cannot adjudicate. We also note the
paper's reported seed sd ($\approx\!0.17$ EM) is far smaller than any arm
here ($0.80$--$3.95$); replications at single seeds in this setting should
be read accordingly.

Limitations. Three seeds: every seed-level interval includes zero. One
model at one scale. BC+ carries $\approx\!\pm2$ EM per-arm eval noise at
830 questions; best-in-window used small subsets; GRPO arms ran one seed;
eval is single-sample greedy; our tier~1 average (four sets) is composed
differently from the paper's MHQA average, so levels are not directly
comparable (our base-model harness anchor matches theirs: $30.8$ vs.\
$31.61$). Nothing here speaks to TreeOPSD at larger scale, with a live
retriever, or on horizons beyond our design.

\section{Conclusion}

Replicated against its own control, TreeOPSD passes. Replicated against a
control that is merely clipped, its long-horizon advantage is zero and
its QA advantage is at most about two points and seed-fragile, while its
mechanism measures at $\approx\!1\%$ of the update it is supposed to
drive. The robust, reusable finding of this study is about the control:
tree backup without a bound collapses, and any bound fixes it. Papers
that compare a bounded method against an unbounded control should expect
this confound to be checked.

\paragraph{Reproducibility.} Preregistration, training and scoring code,
per-run logs and the three seed-package JSONs underlying
Tables~\ref{tab:arms}--\ref{tab:gaps} are released with this note at
\url{https://github.com/sjronan/treeopsd-replication}
(\texttt{PREREG-EXP1.md}, \texttt{scripts/seed\_package.py},
\texttt{results/}). Experiments were
designed and directed by the author and executed with AI-agent assistance
under a file-based controller protocol; all rulings and deviations are
preserved in the released correspondence log.

\begin{thebibliography}{9}
\bibitem{treeopsd} Zou, H., Wang, Q., et al. Tree-Aware Posterior
Distillation for Long-Horizon Agentic Reinforcement Learning.
arXiv:2610.23181, 2026.
\end{thebibliography}

\end{document}

